%% kapa: exam=1 round=36 year=2025 period=1 form=A subject=civil_law subject_ko=민법 seq=8 number=8
%% source: Q-Net 감정평가사 기출문제 (cst00309) artlSeq=5231525 file=raw/1008/5231525/2025년 제36회 감정평가사 1차시험 1교시 문제지 원본.pdf
%% conversion: kordoc 4.24.51 -> tools/gen1.py (auto); review=proofread-sonnet; flags=none
\begin{question}{8}
사단법인에 관한 설명으로 옳지 않은 것은? (다툼이 있으면 판례에 따름)
\begin{choices}
  \item 사단법인의 정관은 특별한 사정이 없는 한 총사원의 3분의 2 이상의 동의가 있는 때에 한하여 변경할 수 있다.
  \item 총사원의 5분의 1 이상이 회의의 목적사항을 제시하여 임시총회의 소집을 청구한 경우, 이사는 특별한 사정이 없는 한 임시총회를 소집해야 한다.
  \item 사단법인은 특별한 사정이 없는 한 총사원의 5분의 3이 동의한 경우에 해산을 결의할 수 있다.
  \item 총회 소집 통지에 목적사항으로 기재되지 않은 사항에 대한 결의는 구성원 전원이 그 총회에 참석하여 의결한 경우가 아니면 원칙적으로 무효이다.
  \item 비법인사단의 사원의 지위는 규약에 의하여 상속될 수 있다.
\end{choices}
\end{question}
\answer{3}
